% =====================================================================
%  Método 7 — Índices de complexidade sintática   (dimensão: estrutura)
%  Origem: células 29, 30 e 31 do caderno
% =====================================================================
\subsection{Índices de complexidade sintática}
\label{met:complexidade}

Os métodos anteriores investigavam a correção da frase. Este investiga outra propriedade, o
grau de elaboração da construção, e o resultado que produz não é um veredito, mas um perfil
numérico. Três índices se extraem da mesma árvore.

O primeiro é a profundidade máxima, definida como a distância do nó mais afastado até a raiz,
%
\begin{equation}
\prof(\sent) = \max_{t \in \sent} \; d(t, \raiz)
\end{equation}
%
em que $d(t,\raiz)$ designa o número de arcos percorridos de $t$ até a raiz. O índice mensura
encaixamento, isto é, o número de níveis de subordinação acumulados pela frase.

O segundo é o número de orações, obtido pela contagem dos nós que encabeçam oração, a
principal e as subordinadas:
%
\begin{equation}
\orac(\sent) = \bigl|\{\, t \in \sent : \mathrm{dep}(t) \in \mathcal{C} \,\}\bigr|,
\qquad \mathcal{C} = \{\texttt{ROOT}, \texttt{ccomp}, \texttt{xcomp}, \texttt{advcl},
\texttt{acl}, \texttt{acl:relcl}, \texttt{csubj}\}
\end{equation}

O terceiro é a ramificação média, o número médio de filhos por nó:
%
\begin{equation}
\ramif(\sent) = \frac{1}{n}\sum_{t \in \sent} \bigl|\mathrm{filhos}(t)\bigr|
\end{equation}

A finalidade destes índices é a comparação de registros --- o jurídico e o jornalístico ---,
a estimativa de dificuldade de leitura e a caracterização autoral. Trata-se de medidas de
estilo e de carga de processamento, e não de correção: uma frase de elevada complexidade não
é por isso incorreta, e uma frase incorreta pode ser trivialmente simples.

O corpus contém duas frases de exatamente 13 palavras e estruturas maximamente distintas. A
primeira constitui lista plana, com um verbo e cinco irmãos dele dependentes; a segunda
encaixa duas relativas, uma no interior da outra. A Tabela~\ref{tab:complexidade} apresenta
os três índices sobre ambas.

\begin{table}[H]
\centering
\dimtable{
\begin{tabular}{@{}lrrrr@{}}
\toprule
\textbf{Frase} & $n$ & $\prof$ & $\orac$ & $\ramif$ \\
\midrule
O incêndio devastou mata, cerrado, pasto, encosta e vale. & 13 & 3 & 1 & 0,9231 \\
O incêndio que atingiu a mata que cercava o vale enfim cessou. & 13 & 6 & 3 & 0,9231 \\
\bottomrule
\end{tabular}}
\caption{Os três índices sobre duas frases de mesmo comprimento.}
\label{tab:complexidade}
\end{table}

A profundidade e o número de orações separam as duas frases com margem ampla: o encaixamento
duplica e a contagem de orações triplica, sem alteração do comprimento. É o comportamento
esperado de um índice de complexidade, a saber, capturar a forma da árvore independentemente
do número de palavras que a compõem, e nenhum dos métodos anteriores distingue essas duas
frases, uma vez que ambas são corretas. A terceira coluna, contudo, é idêntica nas duas, e a
Seção~\ref{sec:empiricos} examina a razão disso.
